% MT015 / explanation — canonical editable content.
NB的“朴素”是：\textbf{已经知道类别后}，把不同特征分别建模。例如在同一类邮件中，近似把各词的出现情况分开估计。它没有保证真实数据一定满足这个假设。



\begin{itemize}

\item \textbf{A按原解选，但要带条件。} 错分代价相同时，选择真实后验最大的类能使出错概率最小。若真实后验为\((0.8,0.2)\)，选第一类错的概率是0.2，选第二类则是0.8。准确估计后验可以保证选中概率最大的类；即使估成\((0.6,0.4)\)，仍选第一类，错误概率也还是0.2。因此“概率估计完全准确”是充分条件，关键是能否选中真实后验最大的类；NB的独立假设本身不保证这一点。

\item \textbf{B不选。} “单个类的分布轮廓”和“两类打平的位置”不是同一幅图。各特征独立使高斯轮廓不倾斜，但比较两类仍可以得到斜边界。

\item \textbf{C选。} NB先估计各类特征分布和先验，再代入Bayes公式算后验，没有直接把类别后验作为训练模型。

\item \textbf{D不选。} 小样本会使均值、方差和词概率估计不稳；碰巧没见过某个词就给它零概率，也是过度贴合训练样本的表现。可以平滑或稳定方差估计。

\item \textbf{E选。} 如果把完全不符合数据的分布强加进去，算出的后验就可能错。因此类条件分布（class-conditional distribution，CCD）的选择会影响分类。

\end{itemize}


看B的一个补充反例：两类先验相同，均值分别为\((0,0)\)、\((1,1)\)，各维方差都是1。边界上到两中心的平方距离相等：



\[
x_1^2+x_2^2=(x_1-1)^2+(x_2-1)^2.
\]


展开右边并抵消平方项，得到\(x_1+x_2=1\)。这是一条斜线，已经足以反驳“总与坐标轴对齐”。


\textbf{出处：} 2021A期中 Q2，原卷第2页（5分）


\textbf{答案依据：} 2021A配套解答 / 2021A配套解答副本 Q2，随卷答案ACE；最优性加成立条件。
